% -----------------------------------------------------------------------------
% BOILERPLATE PRESENTATION LATEX / BEAMER
% Presentation resumee d'un cahier des charges
% Projet : \ph{Nom du projet}
% Client : \ph{Nom du client}
% Auteur : \ph{Nom de l'auteur}
% -----------------------------------------------------------------------------

\documentclass[aspectratio=169,11pt]{beamer}

% -----------------------------------------------------------------------------
% LANGUE ET ENCODAGE
% -----------------------------------------------------------------------------
\usepackage[utf8]{inputenc}
\usepackage[T1]{fontenc}
\usepackage[french]{babel}
\usepackage{lmodern}
\usepackage{microtype}

% -----------------------------------------------------------------------------
% PACKAGES VISUELS ET STRUCTURELS
% -----------------------------------------------------------------------------
\usepackage{graphicx}
\usepackage{xcolor}
\usepackage{booktabs}
\usepackage{tabularx}
\usepackage{array}
\usepackage{multirow}
\usepackage{tikz}
\usepackage{listings}
\usepackage{hyperref}
\usepackage{ragged2e}

% -----------------------------------------------------------------------------
% VARIABLES PRINCIPALES DU DOCUMENT
% -----------------------------------------------------------------------------
\newcommand{\ProjectName}{[Nom du projet]}
\newcommand{\ProjectSubtitle}{Presentation resumee du cahier des charges}
\newcommand{\ClientName}{[Nom du client]}
\newcommand{\AuthorName}{\ph{Nom de l'auteur}}
\newcommand{\TrainingName}{[Formation / organisme]}
\newcommand{\PresentationDate}{[Date]}
\newcommand{\ProjectWebsite}{[URL du site ou depot]}

% -----------------------------------------------------------------------------
% THEME BEAMER
% -----------------------------------------------------------------------------
\usetheme{Madrid}
\usecolortheme{default}
\setbeamertemplate{navigation symbols}{}
\setbeamertemplate{caption}[numbered]
\setbeamertemplate{footline}[frame number]

% Palette modifiable
\definecolor{PrimaryColor}{HTML}{3E2F5B}
\definecolor{SecondaryColor}{HTML}{7A5C99}
\definecolor{AccentColor}{HTML}{C8A86B}
\definecolor{SoftBackground}{HTML}{F7F4EF}
\definecolor{DarkText}{HTML}{1E1E1E}
\definecolor{MutedText}{HTML}{666666}
\definecolor{CodeBackground}{HTML}{F2F2F2}

\setbeamercolor{title}{fg=white,bg=PrimaryColor}
\setbeamercolor{frametitle}{fg=white,bg=PrimaryColor}
\setbeamercolor{structure}{fg=PrimaryColor}
\setbeamercolor{block title}{fg=white,bg=PrimaryColor}
\setbeamercolor{block body}{fg=DarkText,bg=SoftBackground}
\setbeamercolor{alerted text}{fg=AccentColor}

\setbeamerfont{title}{size=\Large,series=\bfseries}
\setbeamerfont{subtitle}{size=\large}
\setbeamerfont{frametitle}{size=\large,series=\bfseries}
\setbeamerfont{block title}{series=\bfseries}

% -----------------------------------------------------------------------------
% CONFIGURATION CODE
% -----------------------------------------------------------------------------
\lstdefinestyle{presentationcode}{
  backgroundcolor=\color{CodeBackground},
  basicstyle=\ttfamily\scriptsize,
  keywordstyle=\bfseries\color{PrimaryColor},
  commentstyle=\itshape\color{MutedText},
  stringstyle=\color{SecondaryColor},
  numbers=left,
  numberstyle=\tiny\color{MutedText},
  stepnumber=1,
  frame=single,
  rulecolor=\color{SecondaryColor},
  breaklines=true,
  showstringspaces=false,
  tabsize=2
}
\lstset{style=presentationcode}

% -----------------------------------------------------------------------------
% COMPOSANTS REUTILISABLES
% -----------------------------------------------------------------------------

% Slide separateur de section
\newcommand{\sectionframe}[2]{%
  \begin{frame}[plain]
    \begin{tikzpicture}[remember picture,overlay]
      \fill[PrimaryColor] (current page.south west) rectangle (current page.north east);
      \fill[AccentColor] (current page.south west) rectangle ([xshift=0.35cm]current page.north west);
    \end{tikzpicture}
    \vfill
    {\color{white}\Huge\bfseries #1}\par\vspace{0.35cm}
    {\color{white}\large #2}
    \vfill
  \end{frame}
}

% Bloc placeholder visuel
\newcommand{\visualplaceholder}[2]{%
  \begin{center}
    \begin{tikzpicture}
      \draw[rounded corners=4pt, dashed, thick, color=SecondaryColor, fill=SoftBackground]
        (0,0) rectangle (#1,#2);
      \node[align=center, color=MutedText] at (#1/2,#2/2) {Visuel a inserer\\[0.1cm]\scriptsize image, schema, capture, mockup};
    \end{tikzpicture}
  \end{center}
}

% Image avec legende courte
\newcommand{\slideimage}[3][0.85\linewidth]{%
  \begin{figure}
    \centering
    \includegraphics[width=#1]{#2}
    \caption{#3}
  \end{figure}
}

% Image pleine largeur avec titre de contexte
\newcommand{\fullvisualslide}[3]{%
  \begin{frame}{#1}
    \centering
    \includegraphics[height=0.72\textheight,width=0.95\linewidth,keepaspectratio]{#2}
    \vspace{0.2cm}

    {\small #3}
  \end{frame}
}

% Deux visuels cote a cote
\newcommand{\twovisuals}[4]{%
  \begin{columns}[T,onlytextwidth]
    \begin{column}{0.49\textwidth}
      \centering
      \includegraphics[width=\linewidth,height=0.48\textheight,keepaspectratio]{#1}
      \scriptsize #2
    \end{column}
    \begin{column}{0.49\textwidth}
      \centering
      \includegraphics[width=\linewidth,height=0.48\textheight,keepaspectratio]{#3}
      \scriptsize #4
    \end{column}
  \end{columns}
}

% Carte KPI simple
\newcommand{\kpicard}[3]{%
  \begin{tikzpicture}
    \node[rounded corners=6pt, fill=SoftBackground, draw=SecondaryColor, line width=0.6pt,
      text width=#1, minimum height=1.7cm, align=center, inner sep=0.15cm] {
      {\Large\bfseries\color{PrimaryColor} #2}\\[-0.05cm]
      {\scriptsize\color{MutedText} #3}
    };
  \end{tikzpicture}
}

% Ligne d'etape pour processus
\newcommand{\processstep}[3]{%
  \begin{tikzpicture}[baseline=(txt.base)]
    \node[circle, fill=PrimaryColor, text=white, minimum size=0.6cm] (num) {#1};
    \node[anchor=west, text width=9.8cm] (txt) at ([xshift=0.25cm]num.east) {\textbf{#2} -- #3};
  \end{tikzpicture}
}

% Bloc court type message strategique
\newcommand{\keymessage}[1]{%
  \begin{block}{Message cle}
    \justifying #1
  \end{block}
}

% Slide vierge pret a remplir
\newcommand{\blankcontentslide}[1]{%
  \begin{frame}{#1}
    \vspace{0.4cm}
    \visualplaceholder{11.5}{5.3}
  \end{frame}
}

% Slide avec notes orales
\newcommand{\speakerhint}[1]{%
  \vspace{0.2cm}
  {\scriptsize\color{MutedText}\textit{Note orale : #1}}
}


% Placeholder texte reutilisable. Utiliser \ph{...} au lieu de crochets en debut de ligne.
\newcommand{\ph}[1]{{\color{MutedText}\texttt{<#1>}}}

% -----------------------------------------------------------------------------
% METADONNEES
% -----------------------------------------------------------------------------
\title{\ProjectName}
\subtitle{\ProjectSubtitle}
\author{\AuthorName}
\institute{\ClientName \\ \TrainingName}
\date{\PresentationDate}

% -----------------------------------------------------------------------------
% DOCUMENT
% -----------------------------------------------------------------------------
\begin{document}

% -----------------------------------------------------------------------------
% COUVERTURE
% -----------------------------------------------------------------------------
\begin{frame}[plain]
  \begin{tikzpicture}[remember picture,overlay]
    \fill[PrimaryColor] (current page.south west) rectangle (current page.north east);
    \fill[AccentColor] (current page.south west) rectangle ([xshift=0.35cm]current page.north west);
  \end{tikzpicture}
  \vfill
  {\color{white}\Huge\bfseries \ProjectName}\par\vspace{0.35cm}
  {\color{white}\Large \ProjectSubtitle}\par\vspace{0.8cm}
  {\color{white}\large Client : \ClientName}\par
  {\color{white}\large Auteur : \AuthorName}\par
  {\color{white}\large Date : \PresentationDate}
  \vfill
  {\color{white}\small \ProjectWebsite}
\end{frame}

% -----------------------------------------------------------------------------
% OBJECTIF DE LA PRESENTATION
% -----------------------------------------------------------------------------
\begin{frame}{Objectif de la presentation}
  \begin{columns}[T,onlytextwidth]
    \begin{column}{0.58\textwidth}
      \keymessage{Cette presentation resume les points principaux du cahier des charges afin de donner une vision claire du projet, de ses objectifs, de ses contraintes et de son deroulement.}
      \begin{itemize}
        \item Clarifier le contexte du projet.
        \item Presenter les besoins principaux.
        \item Montrer la cible, les fonctions et les contraintes.
        \item Resumer la planification, les livrables et le budget.
      \end{itemize}
    \end{column}
    \begin{column}{0.38\textwidth}
      \visualplaceholder{4.7}{4.2}
    \end{column}
  \end{columns}
\end{frame}

\begin{frame}{Sommaire}
  \tableofcontents
\end{frame}

% -----------------------------------------------------------------------------
% INTRODUCTION
% -----------------------------------------------------------------------------
\section{Introduction}
\sectionframe{Introduction}{Presenter le projet en quelques phrases}

\begin{frame}{Introduction du projet}
  \begin{columns}[T,onlytextwidth]
    \begin{column}{0.55\textwidth}
      \begin{block}{Resume}
        \ph{Decrire le projet en 3 a 5 phrases.}
      \end{block}
      \begin{block}{Problematique}
        \ph{Problematique a completer.}
      \end{block}
    \end{column}
    \begin{column}{0.41\textwidth}
      \visualplaceholder{5.0}{4.6}
    \end{column}
  \end{columns}
\end{frame}

\begin{frame}{Vision synthetique}
  \begin{center}
    \kpicard{3.0cm}{\ph{1}}{Activite principale}\hspace{0.35cm}
    \kpicard{3.0cm}{\ph{2}}{Publics cibles}\hspace{0.35cm}
    \kpicard{3.0cm}{\ph{3}}{Offres / services}
  \end{center}
  \vspace{0.5cm}
  \begin{block}{A retenir}
    \ph{Phrase de synthese}
  \end{block}
\end{frame}

% -----------------------------------------------------------------------------
% PRESENTATION DETAILLEE
% -----------------------------------------------------------------------------
\section{Presentation detaillee}
\sectionframe{Presentation detaillee}{Decrire l'activite, les offres et l'identite du projet}

\begin{frame}{Presentation de l'activite}
  \begin{columns}[T,onlytextwidth]
    \begin{column}{0.5\textwidth}
      \begin{itemize}
        \item Nom commercial : \ph{Nom}.
        \item Domaine : \ph{Domaine}.
        \item Format : \ph{Format}.
        \item Ton : \ph{Ton}.
      \end{itemize}
    \end{column}
    \begin{column}{0.46\textwidth}
      \visualplaceholder{5.7}{4.0}
    \end{column}
  \end{columns}
\end{frame}

\begin{frame}{Offres principales}
  \begin{tabularx}{\textwidth}{>{\bfseries}p{3.2cm}X p{2.5cm}}
    \toprule
    Offre & Description courte & Format \\
    \midrule
    \ph{Offre 1} & \ph{Description synthetique} & \ph{En ligne} \\
    \ph{Offre 2} & \ph{Description synthetique} & \ph{Collectif} \\
    \ph{Offre 3} & \ph{Description synthetique} & \ph{Presentiel} \\
    \ph{Offre 4} & \ph{Description synthetique} & \ph{Sur mesure} \\
    \bottomrule
  \end{tabularx}
\end{frame}

\begin{frame}{Identite et univers visuel}
  \begin{columns}[T,onlytextwidth]
    \begin{column}{0.33\textwidth}
      \begin{block}{Ambiance}
        \ph{Mots-cles}
      \end{block}
    \end{column}
    \begin{column}{0.33\textwidth}
      \begin{block}{Couleurs}
        \ph{Palette}
      \end{block}
    \end{column}
    \begin{column}{0.33\textwidth}
      \begin{block}{Typographies}
        \ph{Typographies}
      \end{block}
    \end{column}
  \end{columns}
  \vspace{0.3cm}
  \visualplaceholder{11.5}{2.2}
\end{frame}

% -----------------------------------------------------------------------------
% INTERVENANTS
% -----------------------------------------------------------------------------
\section{Intervenants}
\sectionframe{Intervenants}{Identifier les roles et responsabilites}

\begin{frame}{Parties prenantes}
  \begin{tabularx}{\textwidth}{p{3.2cm}p{3.0cm}X}
    \toprule
    Nom / role & Responsabilite & Besoins d'information \\
    \midrule
    \ph{Client} & \ph{Responsabilites} & \ph{Informations} \\
    \ph{Developpeur} & \ph{Responsabilites} & \ph{Documents} \\
    \ph{Referent} & \ph{Responsabilites} & \ph{Suivi} \\
    \ph{Utilisateur final} & \ph{Usage} & \ph{Attentes} \\
    \bottomrule
  \end{tabularx}
\end{frame}

\begin{frame}{Organisation du projet}
  \processstep{1}{Recueil du besoin}{Collecte des objectifs, contenus, contraintes et inspirations.}\vspace{0.25cm}

  \processstep{2}{Conception}{Definition de l'arborescence, des pages, des parcours et des maquettes.}\vspace{0.25cm}

  \processstep{3}{Realisation}{Integration, developpement, contenu, responsive design et tests.}\vspace{0.25cm}

  \processstep{4}{Validation}{Recette, corrections, livraison et mise en ligne.}
\end{frame}

% -----------------------------------------------------------------------------
% CONTEXTE
% -----------------------------------------------------------------------------
\section{Contexte}
\sectionframe{Contexte}{Objectifs, cible, concurrence et analyse marketing}

\begin{frame}{Objectifs du projet}
  \begin{columns}[T,onlytextwidth]
    \begin{column}{0.48\textwidth}
      \begin{block}{Objectifs principaux}
        \begin{itemize}
          \item \ph{Objectif principal 1}
          \item \ph{Objectif principal 2}
          \item \ph{Objectif principal 3}
        \end{itemize}
      \end{block}
    \end{column}
    \begin{column}{0.48\textwidth}
      \begin{block}{Indicateurs de reussite}
        \begin{itemize}
          \item \ph{Nombre de contacts}
          \item \ph{Taux de conversion}
          \item \ph{Lisibilite}
          \item \ph{Prise de contact}
        \end{itemize}
      \end{block}
    \end{column}
  \end{columns}
\end{frame}

\begin{frame}{Cible utilisateur}
  \begin{columns}[T,onlytextwidth]
    \begin{column}{0.4\textwidth}
      \visualplaceholder{4.8}{4.7}
    \end{column}
    \begin{column}{0.56\textwidth}
      \begin{block}{Persona principal}
        \textbf{Nom :} \ph{Persona}\newline
        \textbf{Age :} \ph{Tranche age}\newline
        \textbf{Besoin :} \ph{Besoin}\newline
        \textbf{Frein :} \ph{Frein}\newline
        \textbf{Attente :} \ph{Attente}
      \end{block}
    \end{column}
  \end{columns}
\end{frame}

\begin{frame}{Concurrence}
  \begin{tabularx}{\textwidth}{p{3.0cm}X X X}
    \toprule
    Concurrent & Points forts & Points faibles & Opportunite \\
    \midrule
    \ph{Concurrent 1} & \ph{A completer} & \ph{A completer} & \ph{A completer} \\
    \ph{Concurrent 2} & \ph{A completer} & \ph{A completer} & \ph{A completer} \\
    \ph{Concurrent 3} & \ph{A completer} & \ph{A completer} & \ph{A completer} \\
    \bottomrule
  \end{tabularx}
\end{frame}

\begin{frame}{Analyse marketing resumee}
  \begin{columns}[T,onlytextwidth]
    \begin{column}{0.48\textwidth}
      \begin{block}{Positionnement}
        \ph{Positionnement}
      \end{block}
      \begin{block}{Message commercial}
        \ph{Promesse}
      \end{block}
    \end{column}
    \begin{column}{0.48\textwidth}
      \begin{block}{Canaux possibles}
        \begin{itemize}
          \item Site web vitrine.
          \item Reseaux sociaux.
          \item Referencement naturel.
          \item Recommandations / bouche-a-oreille.
        \end{itemize}
      \end{block}
    \end{column}
  \end{columns}
\end{frame}

% -----------------------------------------------------------------------------
% DEFINITION DES BESOINS
% -----------------------------------------------------------------------------
\section{Definition des besoins}
\sectionframe{Definition des besoins}{Existant, besoin, fonctions du produit}

\begin{frame}{Etude de l'existant}
  \begin{columns}[T,onlytextwidth]
    \begin{column}{0.52\textwidth}
      \begin{block}{Constat}
        \ph{Etat de l'existant}
      \end{block}
      \begin{block}{Limites observees}
        \begin{itemize}
          \item \ph{Limite 1}
          \item \ph{Limite 2}
          \item \ph{Limite 3}
        \end{itemize}
      \end{block}
    \end{column}
    \begin{column}{0.44\textwidth}
      \visualplaceholder{5.4}{4.5}
    \end{column}
  \end{columns}
\end{frame}

\begin{frame}{Enonce du besoin}
  \begin{block}{Besoin principal}
    \ph{Enonce du besoin}
  \end{block}
  \begin{block}{Synthese utilisateur}
    En tant que \textbf{\ph{type d'utilisateur}}, je veux \textbf{\ph{action ou information}}, afin de \textbf{\ph{benefice attendu}}.
  \end{block}
  \speakerhint{Presenter le besoin avec des mots simples, sans entrer dans les details techniques.}
\end{frame}

\begin{frame}{Fonctions principales du produit}
  \begin{tabularx}{\textwidth}{p{1.7cm}X p{2.5cm}p{2.2cm}}
    \toprule
    ID & Fonction & Priorite & Statut \\
    \midrule
    F-01 & Presenter l'activite et son univers. & Essentielle & \ph{A faire} \\
    F-02 & Afficher les offres et leurs formats. & Essentielle & \ph{A faire} \\
    F-03 & Permettre la prise de contact. & Essentielle & \ph{A faire} \\
    F-04 & Rassurer avec une presentation claire. & Importante & \ph{A faire} \\
    F-05 & Optimiser le site pour mobile. & Essentielle & \ph{A faire} \\
    \bottomrule
  \end{tabularx}
\end{frame}

\begin{frame}{Parcours utilisateur resume}
  \begin{center}
    \begin{tikzpicture}[node distance=1.7cm, every node/.style={align=center}]
      \node[draw, rounded corners, fill=SoftBackground, minimum width=2.7cm, minimum height=1cm] (a) {Arrivee\\sur le site};
      \node[draw, rounded corners, fill=SoftBackground, right of=a, xshift=2.1cm, minimum width=2.7cm, minimum height=1cm] (b) {Lecture\\des offres};
      \node[draw, rounded corners, fill=SoftBackground, right of=b, xshift=2.1cm, minimum width=2.7cm, minimum height=1cm] (c) {Choix\\d'une seance};
      \node[draw, rounded corners, fill=SoftBackground, right of=c, xshift=2.1cm, minimum width=2.7cm, minimum height=1cm] (d) {Contact\\ou reservation};
      \draw[->, thick, color=PrimaryColor] (a) -- (b);
      \draw[->, thick, color=PrimaryColor] (b) -- (c);
      \draw[->, thick, color=PrimaryColor] (c) -- (d);
    \end{tikzpicture}
  \end{center}
  \vspace{0.4cm}
  \begin{block}{Point d'attention}
    \ph{Point d'attention}
  \end{block}
\end{frame}

% -----------------------------------------------------------------------------
% CONTRAINTES
% -----------------------------------------------------------------------------
\section{Contraintes}
\sectionframe{Contraintes}{Contraintes techniques, legales, budgetaires et temporelles}

\begin{frame}{Contraintes techniques}
  \begin{columns}[T,onlytextwidth]
    \begin{column}{0.48\textwidth}
      \begin{block}{Environnement}
        \begin{itemize}
          \item \ph{Framework / CMS / technologie}
          \item \ph{Hebergement}
          \item \ph{Nom de domaine}
          \item \ph{Maintenance}
        \end{itemize}
      \end{block}
    \end{column}
    \begin{column}{0.48\textwidth}
      \begin{block}{Qualite attendue}
        \begin{itemize}
          \item Responsive design.
          \item Performance correcte.
          \item Accessibilite de base.
          \item Sauvegarde et securite minimales.
        \end{itemize}
      \end{block}
    \end{column}
  \end{columns}
\end{frame}

\begin{frame}[fragile]{Slide technique : exemple de pile projet}
  \begin{lstlisting}[language=bash]
# Exemple a adapter
Frontend : [Next.js / React / HTML-CSS]
Style    : [Tailwind / CSS modules / autre]
Backend  : [Aucun / API / CMS / autre]
Hosting  : [OVH / Infomaniak / Vercel / serveur VPS]
Domaine  : [nomdedomaine.fr]
  \end{lstlisting}
  \speakerhint{Utiliser ce format pour presenter rapidement une information technique sans surcharger la slide.}
\end{frame}

\begin{frame}{Contraintes legales et reglementaires}
  \begin{tabularx}{\textwidth}{p{3.2cm}X p{2.4cm}}
    \toprule
    Sujet & Application au projet & Priorite \\
    \midrule
    RGPD & \ph{Formulaire, donnees personnelles, consentement} & Haute \\
    Mentions legales & \ph{Identite editeur, hebergeur, contact} & Haute \\
    Cookies & \ph{Analyse, audience, consentement} & Moyenne \\
    Accessibilite & \ph{Lisibilite, contraste, navigation} & Moyenne \\
    Droit d'auteur & \ph{Images, textes, polices, licences} & Haute \\
    \bottomrule
  \end{tabularx}
\end{frame}

\begin{frame}{Contraintes de couts}
  \begin{center}
    \kpicard{3.2cm}{\ph{XXX EUR}}{Budget initial}\hspace{0.5cm}
    \kpicard{3.2cm}{\ph{XX EUR/mois}}{Couts recurrents}\hspace{0.5cm}
    \kpicard{3.2cm}{\ph{X jours}}{Temps estime}
  \end{center}
  \vspace{0.5cm}
  \begin{block}{Commentaires}
    \ph{Commentaires budget}
  \end{block}
\end{frame}

\begin{frame}{Contraintes de delais}
  \begin{tabularx}{\textwidth}{p{3.0cm}p{3.0cm}X}
    \toprule
    Phase & Date cible & Resultat attendu \\
    \midrule
    Cadrage & \ph{Date} & Cahier des charges valide \\
    Maquettes & \ph{Date} & Structure et direction visuelle validees \\
    Developpement & \ph{Date} & Version fonctionnelle \\
    Recette & \ph{Date} & Corrections et tests \\
    Livraison & \ph{Date} & Mise en ligne ou rendu final \\
    \bottomrule
  \end{tabularx}
\end{frame}

% -----------------------------------------------------------------------------
% DEROULEMENT DU PROJET
% -----------------------------------------------------------------------------
\section{Deroulement du projet}
\sectionframe{Deroulement du projet}{Livrables, planning, devis et validation}

\begin{frame}{Livrables attendus}
  \begin{columns}[T,onlytextwidth]
    \begin{column}{0.48\textwidth}
      \begin{block}{Documents}
        \begin{itemize}
          \item Cahier des charges.
          \item Arborescence.
          \item Maquettes ou wireframes.
          \item Documentation technique simple.
        \end{itemize}
      \end{block}
    \end{column}
    \begin{column}{0.48\textwidth}
      \begin{block}{Produit}
        \begin{itemize}
          \item Site ou prototype fonctionnel.
          \item Pages principales integrees.
          \item Formulaire ou moyen de contact.
          \item Version responsive.
        \end{itemize}
      \end{block}
    \end{column}
  \end{columns}
\end{frame}

\begin{frame}{Planning simplifie}
  \begin{center}
    \begin{tikzpicture}
      \draw[->, thick, color=PrimaryColor] (0,0) -- (11.5,0);
      \foreach \x/\label in {0.5/Cadrage,2.8/Conception,5.1/Integration,7.4/Tests,9.7/Livraison} {
        \draw[fill=AccentColor, draw=PrimaryColor] (\x,0) circle (0.13);
        \node[align=center, above] at (\x,0.15) {\small \label};
        \node[align=center, below] at (\x,-0.15) {\scriptsize \ph{Date}};
      }
    \end{tikzpicture}
  \end{center}
  \vspace{0.5cm}
  \begin{block}{Hypothese}
    \ph{Hypothese planning}
  \end{block}
\end{frame}

\begin{frame}{Devis resume}
  \begin{tabularx}{\textwidth}{X p{2.2cm}p{2.2cm}p{2.2cm}}
    \toprule
    Poste & Quantite & Prix unitaire & Total \\
    \midrule
    Cadrage et conception & \ph{X} & \ph{EUR} & \ph{EUR} \\
    Maquettes & \ph{X} & \ph{EUR} & \ph{EUR} \\
    Integration / developpement & \ph{X} & \ph{EUR} & \ph{EUR} \\
    Tests et corrections & \ph{X} & \ph{EUR} & \ph{EUR} \\
    Mise en ligne & \ph{X} & \ph{EUR} & \ph{EUR} \\
    \midrule
    \textbf{Total estime} & & & \textbf{\ph{EUR}} \\
    \bottomrule
  \end{tabularx}
\end{frame}

\begin{frame}{Risques et points de vigilance}
  \begin{tabularx}{\textwidth}{p{3.2cm}X p{2.3cm}}
    \toprule
    Risque & Impact possible & Prevention \\
    \midrule
    Contenus incomplets & Retard de production & Liste de contenus a fournir \\
    Perimetre instable & Derive du planning & Validation par phase \\
    Contraintes legales & Corrections tardives & Verification RGPD / mentions \\
    Visuels indisponibles & Site moins qualitatif & Banque d'images ou creation graphique \\
    \bottomrule
  \end{tabularx}
\end{frame}

% -----------------------------------------------------------------------------
% SUPPORTS VISUELS
% -----------------------------------------------------------------------------
\section{Supports visuels}
\sectionframe{Supports visuels}{Slides pretes pour captures, maquettes, schemas et images}

\blankcontentslide{Capture d'ecran a commenter}

\begin{frame}{Comparaison avant / apres}
  \begin{columns}[T,onlytextwidth]
    \begin{column}{0.48\textwidth}
      \centering
      \visualplaceholder{5.5}{4.2}
      \textbf{Avant}\newline
      \scriptsize \ph{Etat initial}
    \end{column}
    \begin{column}{0.48\textwidth}
      \centering
      \visualplaceholder{5.5}{4.2}
      \textbf{Apres}\newline
      \scriptsize \ph{Solution proposee}
    \end{column}
  \end{columns}
\end{frame}

\begin{frame}{Mockup ou wireframe}
  \visualplaceholder{11.8}{5.2}
  \speakerhint{Utiliser cette slide pour montrer une maquette pleine largeur.}
\end{frame}

\begin{frame}{Schema technique simplifie}
  \begin{center}
    \begin{tikzpicture}[node distance=2cm]
      \node[draw, rounded corners, fill=SoftBackground, minimum width=2.6cm, minimum height=1cm] (u) {Utilisateur};
      \node[draw, rounded corners, fill=SoftBackground, right of=u, xshift=2.5cm, minimum width=2.6cm, minimum height=1cm] (s) {Site web};
      \node[draw, rounded corners, fill=SoftBackground, right of=s, xshift=2.5cm, minimum width=2.6cm, minimum height=1cm] (h) {Hebergement};
      \draw[->, thick, color=PrimaryColor] (u) -- (s);
      \draw[->, thick, color=PrimaryColor] (s) -- (h);
    \end{tikzpicture}
  \end{center}
  \begin{block}{Explication}
    \ph{Explication}
  \end{block}
\end{frame}

\begin{frame}[fragile]{Exemple de slide technique avec code}
  \begin{columns}[T,onlytextwidth]
    \begin{column}{0.56\textwidth}
      \begin{lstlisting}[language=Java]
const service = {
  name: "\ph{Nom du service}",
  format: "\ph{en ligne / presentiel}",
  active: true,
};
      \end{lstlisting}
    \end{column}
    \begin{column}{0.4\textwidth}
      \begin{block}{Commentaire}
        \ph{Commentaire}
      \end{block}
      \visualplaceholder{4.7}{2.2}
    \end{column}
  \end{columns}
\end{frame}

% -----------------------------------------------------------------------------
% SYNTHESE ET CONCLUSION
% -----------------------------------------------------------------------------
\section{Synthese}
\sectionframe{Synthese}{Conclusion et decision attendue}

\begin{frame}{Synthese du cahier des charges}
  \begin{itemize}
    \item Le projet vise a \ph{objectif principal}.
    \item La cible prioritaire est \ph{cible principale}.
    \item Les fonctions essentielles sont \ph{fonctions principales}.
    \item Les contraintes majeures sont \ph{contraintes principales}.
    \item Les livrables attendus sont \ph{livrables principaux}.
  \end{itemize}
  \vspace{0.3cm}
  \keymessage{\ph{Phrase finale}}
\end{frame}

\begin{frame}{Prochaines etapes}
  \processstep{1}{Valider le perimetre}{Confirmer les pages, fonctions et contraintes.}\vspace{0.25cm}

  \processstep{2}{Fournir les contenus}{Textes, images, informations legales, tarifs.}\vspace{0.25cm}

  \processstep{3}{Lancer la production}{Maquettes, integration, tests et corrections.}\vspace{0.25cm}

  \processstep{4}{Livrer}{Recette finale, documentation et mise en ligne.}
\end{frame}

\begin{frame}[plain]
  \begin{tikzpicture}[remember picture,overlay]
    \fill[PrimaryColor] (current page.south west) rectangle (current page.north east);
  \end{tikzpicture}
  \vfill
  \begin{center}
    {\color{white}\Huge\bfseries Merci}\par\vspace{0.4cm}
    {\color{white}\Large Questions / remarques}\par\vspace{0.8cm}
    {\color{white}\small \AuthorName -- \PresentationDate}
  \end{center}
  \vfill
\end{frame}

% -----------------------------------------------------------------------------
% ANNEXES
% -----------------------------------------------------------------------------
\appendix
\section{Annexes}
\sectionframe{Annexes}{Slides supplementaires a utiliser si necessaire}

\begin{frame}{Annexe : sources et references}
  \begin{tabularx}{\textwidth}{p{3.5cm}X p{2.8cm}}
    \toprule
    Source & Utilisation & Statut \\
    \midrule
    \ph{Source 1} & \ph{Utilisation} & \ph{A verifier} \\
    \ph{Source 2} & \ph{Utilisation} & \ph{A verifier} \\
    \ph{Source 3} & \ph{Utilisation} & \ph{A verifier} \\
    \bottomrule
  \end{tabularx}
\end{frame}

\begin{frame}{Annexe : backlog simplifie}
  \begin{tabularx}{\textwidth}{p{1.5cm}X p{2.2cm}p{2.2cm}}
    \toprule
    ID & Tache & Priorite & Statut \\
    \midrule
    T-01 & \ph{Tache a realiser} & Haute & A faire \\
    T-02 & \ph{Tache a realiser} & Moyenne & A faire \\
    T-03 & \ph{Tache a realiser} & Basse & Optionnel \\
    \bottomrule
  \end{tabularx}
\end{frame}

\begin{frame}{Annexe : slide vide}
  \visualplaceholder{11.8}{5.3}
\end{frame}

\end{document}
